\chapter{Ocean Vuong Teaches the Art of Writing}

この章は、LazyLearn writing and speaking collection における Ocean Vuong と host との conversation を追うものであり、LazyingArt LLC によって curated されている。講義は admiration を method へ変換するところから始まる。awe、freshness、enchantment、wonder、これらはすべて practical な craft question になる。どうすれば language にもう一度 see させることができるのか。no validated board or slide mathematics survives for this lecture なので、以下の displayed formulas は visual transcriptions ではなく、spoken structure を transcript に基づいて condense したものである。ここでそれらが useful なのは、この講義自体が、contrasts、transfers、thresholds、causal chains によって繰り返し thinking しているからだ。

\section{metaphor、observation、そして mimesis からの first break}

host は、method に name を与える前に、Vuong の work のある quality に name を与えるところから始める。writing の中には wonder がある。だが、それだけではなく seeing の way の中にも wonder がある。そして host は、そのような seeing は muscle のように cultivate できるのかと ask する。Vuong の最初の answer は immediate で technical である。metaphor こそが、この question が teachable になる place なのだ、と。students が good metaphor をどう write するのかと ask したなら、first answer は brilliance ではないし ornament でもない。observation である。world を長いあいだ look することだ。the rest は arrangement と syntax だ、と彼は言う。

\begin{definition}
tenor を、description の under にある thing とし、vehicle を、その description を elsewhere へ carry する thing とする。すると Vuong の metaphor のための opening schema は
\begin{equation}
(\text{tenor},\text{vehicle}) \Longrightarrow \text{new correspondence} \Longrightarrow \text{re-seeing}.
\end{equation}
となる。metaphor が succeed するのは、その correspondence が reader に tenor を merely recognize させるのではなく、あらためて perceive させるときである。
\end{definition}

だからこそ講義は、ただちに Aristotle 的な二つの sentence-functions の split へ手を伸ばす。
\begin{align}
\text{mimesis} &:\ \text{known scene} \mapsto \text{recognizable description},\\
\text{poiesis} &:\ \text{known scene} \mapsto \text{estranged process}.
\end{align}
mimetic sentence は、すでに自分たちが name する仕方を know しているものへ私たちを back させる。poetic sentence は、その automatic naming を interrupt する。その spirit の中で、Vuong の opening rule は
\begin{equation}
\text{metaphor} \neq \text{ornament}, \qquad \text{metaphor} = \text{observation} + \text{arrangement/syntax}.
\end{equation}
と書ける。

lecture の first worked example は、Isaac Babel の Red Cavalry における sunset である。Vuong は、新聞記事ふうの line、``a red evening sunset along the hills'' と、Babel の low red sun が ``as if beheaded'' に hills を roll していく line とを contrast する。point は、Babel が merely vivid だということではない。historical situation が sentence の内側に carry されていることだ。Babel が war の中から write しているのだということを、external commentary に tell してもらう必要はない。image 自身がその work をしているのである。しかもそれは、scene の felt motion そのものを change する。

\paragraph{Worked Example: Babel's Sunset.}
\begin{align}
\text{sunset} + \text{beheading image} &\Longrightarrow \text{war context embedded in perception},\\
\text{embedded violence} &\Longrightarrow \text{altered rate of the sunset}.
\end{align}
\begin{enumerate}
\item 私たちは、完全に recognizable な scene から begin する。hills を over した sunset である。
\item sunset を merely decorate するのではなく、violence をその inside に carry する image を add する。
\item その added image が、historical consciousness を perception それ自体の中に embed する。
\item sunset はいまや、roll し、fall し、its head を lose するように seem する。second clause が scene の speed を change するのである。
\item したがってこの line は、merely describe する以上のことをする。reader の seeing を reorganize する。
\end{enumerate}

Vuong は、この ambition に deliberately severe な slogan を与える。
\begin{equation}
\text{competent sentence} \neq \text{sentence the species has never had yet}.
\end{equation}
left-hand side が useless だということではない。私たちには still scaffolding sentences が必要である。だが lecture はまず insist する。writer の deeper task は、merely competent recognition を produce することではない。new perceptual event を possible にすることなのだ。

\subsection{Question \& Answer}

どうすれば、world を merely mimic する以上のことをする metaphor が書けるのか。

私たちは cleverness を chase するところから始めない。object の ready-made label が inadequate に feel されるまで look することから始める。そのあと、object の pressure、pace、あるいは emotional weather を change する vehicle を search する。sequence は practical である。
\begin{enumerate}
\item tenor を identify し、それが generic に feel されなくなるまで looking を続ける。
\item stock resemblance からではなく、sustained observation から vehicle を choose する。
\item scene を merely embellish するのではなく、それを alter する correspondence を build する。
\item reader が merely recognize するのではなく see させられるのでなければ、その line を keep しない。
\end{enumerate}
だからこそ strong な metaphor は、arrive するまで years を要することさえある。hard part は syntax alone ではない。right correspondence を earn することなのである。

\section{correction ではなく recognition としての workshop}

ひとたび lecture が live な sentence の doing することを示したあと、host は lens を narrow する。もし writing が timid にも alive にもなりうるなら、writers にどうやって自分自身の work におけるその difference を notice させるのか。Vuong の answer は、workshop-as-correction を workshop-as-recognition へ replace することである。この shift は subtle だが decisive である。workshop は、more comments、more drafts、more process を produce する merely その fact によって automatically help するわけではない。

lecture はこの contrast を、次のように compress する。
\begin{align}
\text{pattern recognition} &\Longrightarrow \text{self-recognition} \Longrightarrow \text{more intelligent revision},\\
\text{dogmatic correction} &\Longrightarrow \text{homogenization or destruction of the work}.
\end{align}
ここが matter するのは、broader culture が私たちを factory のように thinking させたがるからである。work を machine に feed し、more procedure を apply すれば、それは better になるはずだ、と。Vuong はその fantasy を reject する。some workshops は work を destroy する。some drafts は、work の pinnacle を越えてしまう。productivity のように見えるものが rubble になることもあるのである。

lecture の clearest formal definition は、ここで現れる。
\begin{equation}
\text{sentence} = \text{consciousness filtered through syntax}.
\end{equation}
もしこれが right なら、workshop は、draft がすでにどんな consciousness を reveal しているのかを ask することから始まらねばならない。どの images が recur しているか。どの verbs が line breaks をまたいで落ちてくるのか。どこで paragraph が tense を switch するのか。どの prepositions が line を forward に launch し続けているのか。Vuong の classroom examples が deliberately concrete なのは、point が abstract な ``theme'' にないからである。point は、mind がすでに move している formal habits を see することにある。

だからこそ lecture は、overproduction の dangers に return する。poet が一年間毎日一篇 poem を write したからといって、more usable work を end up するとは限らない。fresh beginning より salvage が harder な field of debris を end up するかもしれない。これに対して recognition は、より simpler な question を ask する。ここで何がすでに起こっているのか。ひとたびその question が serious になれば、revision は dogma の application であることを stop し、tendency の clarification になる。

Japanese botanist の anecdote は、これを method に turn する。どうして এত many medicinal plants を見つけられたのかと ask され、その botanist は、自分は known medicine に look するものを search したのではない、simply 自分にとって new なものを search し、それが medicine であることを hope したのだ、と言う。sometimes それは medicine ではなかった。sometimes それは poison だった。だが field を enlarge したのは、approved model に resemblance していることではなく、novelty-seeking だった。Vuong は clearly にこれを writing classroom の model として offer している。

\subsection{Question \& Answer}

なぜ more feedback と more drafting が、piece を better にするどころか worse にしてしまいうるのか。

more process は better recognition と same thing ではないからである。revision は、work 自体がまだ earn していない inherited rules に governed されるとき fail する。Vuong の alternative は procedural だが mechanical ではない。
\begin{enumerate}
\item repair を propose する前に patterns を read する。
\item diction、tense、syntax、image、line-movement における recurring formal tendencies に name を与える。
\item piece が unexpectedly alive になる place を ask する。
\item generic correctness に toward するのではなく、それらの discoveries に toward revise する。
\end{enumerate}
この model では、workshop は poem を manufacture しない。それがすでに trying to become していたものへ writer が sooner に arrive するのを help するのである。

\section{sentence はどう tame されたか}

interruption のあと、lecture は recap によって return しない。tension によって return する。host は novelty、surprise、freshness、そして quality の pursuit が it itself restrictive になってしまうのではないかと ask する。Vuong は answer によって frame を widen する。modern sentence は、彼に言わせれば、timeless standard of good writing ではない。historical product なのである。

その larger history は、prior question から始まる。そもそも literature とは何を mean するのか。Vuong は insist する。category それ自体が late なのである。modern department と、その institutional sorting の before には、poem、story、letter、speech、practical writing の境界はもっと looser だった。novel でさえ、always serious literary instrument として treat されていたわけではない。とりわけ American context において Civil War のあとに起きたその later moral elevation は、もう一つの development、すなわち journalism の standardization と coincide していたのである。

lecture の historical contrast は、schematically には
\begin{align}
\text{oratorical / Victorian sentence} &\Longrightarrow \text{delay} + \text{subordination} + \text{buildup},\\
\text{newspaper sentence} &\Longrightarrow \text{brevity} + \text{clarity} + \text{efficiency} + \text{standardization}.
\end{align}
と書ける。older sentence は sermon、speech、public argument に belong しており、rhythm と suspension が attention を keep alive している culture に belong している。newer sentence は information delivery に belong している。それは reliable で、clear で、fast で、circulation や advertising と coexist するに enough brief でなければならないのである。

だからこそ Churchill の anaphora は、random digression としてではなく、older sentence の live memory として lecture に appear するのである。``We shall fight...'' は、public force を buildup しながら payoff を delay する。Vuong はそれから、partially corrupted な transcript stretch からも preserve できる stable claim へ cautiously だが clearly に return する。twentieth-century prose は newspaper model を absorb し、Hemingway、Crane、London、Orwell のような writers は、私たちの culture がいま often simply ``good writing'' と呼ぶ clipped、efficient sentence の stand-ins になったのだ、と。

argument は、host が right angles の image を offer するとき、more pointed になる。much contemporary prose は、ruler-straight で、angular で、coarse で、over-refined に feel されるのだ、と彼は言う。Vuong はその analogy を seize する。
\begin{equation}
\text{industrial right angle} \Longrightarrow \text{cultural preference for standardized form}.
\end{equation}
これは literal geometry ではない。industrial modernity が、clean edges、repeatability、mass-producible surfaces を value するよう私たちを teach してきた、その historical drift を describe する one way なのである。prose も同じ aspiration を acquire する。sentence は ``invisible'' で、efficient で、authorial pressure を suspicious に feel するようになる。

そこから lecture は naturally に contemporary examples へ move する。Microsoft Word は Shakespeare に correction marks を浴びせる。AI は、human language への sudden betrayal としてではなく、ずっと older な homogenization の project の continuation として appear する。automated prose の long before に、私たちはすでに sentence を sameness へ toward training していたのだ。

\subsection{Question \& Answer}

innovation を praise している even ときに、なぜ contemporary prose は so often standardized に sound するのか。

私たちの institutions が innovation を retrospectively には praise しながら、procedurally には suppress するからである。私たちは daring な masters を teach しながら、living student には、less conspicuous に、less idiosyncratic に、existing product と less unlike に sound するよう ask する。editing、review culture、software、market fear は all recognizability を reward する。result は prose の death ではない。prose が sound してよいとされるものの narrowing である。

\section{estrangement、thresholds、そして again see することを学ぶ}

この narrowing が diagnosed されると、host は disenchantment を directly に name する。Vuong は、middle lecture の key counter-principle で answer する。estrangement である。ここで Viktor Shklovsky が central になる。もっとも bold な claim は、object 自体の中に cliché などというものはない、ということだ。cliché は exhausted な treatment の中に reside する。もし roses、grandmothers、kitchens、mountains、flowers を ``already used'' だとして ban し続けたなら、ultimately 私たちは world そのものを literature から exile してしまうことになる。

basic structure は、
\begin{equation}
\text{familiar object} + \text{unexpected placement} \Longrightarrow \text{fresh perception}.
\end{equation}
となる。したがって lecture の sharpest illustration は、
\begin{equation}
\text{rose in bridesmaid's hair} \neq \text{rose in Mike Tyson's ear}.
\end{equation}
である。rose 自体は unchanged である。treatment が changed したのである。したがって lesson は ``never write the rose'' ではない。``reconsider the rose'' なのである。

Shklovsky が quote する Tolstoy は、その point を deepen する。Tolstoy は、自分が sofa の dusting をしたかどうか remember できない。それが routine で unconscious な act だったからである。もし life が routine のみのうちに lived されるなら、それは almost、never really lived されたことがないかのように pass してしまう。literature がここで matter するのは、routine を interrupt するからである。それは、seeing と merely recognizing との difference を restore するのである。

host の analogies は、lecture のこの portion において especially important である。抽象を concrete に keep してくれるからだ。Bierstadt は mountain を visible again にする。Monet の lilies と Van Gogh の flowers は、mere な objecthood だけでなく energy を show する。それから host は、video editing と early photography へ手を伸ばす。frames に zoom into してみると、single smooth act に見えたものが、multiple distinct moments として reveal されるのだ。これは exactly、Vuong が perception に望んでいることなのである。writing は、world を slow し、its intermediate states と encounter できるようにしなければならない。

これによって lecture は Aristotle へ return する。
\begin{equation}
\text{bud} \leadsto \text{rose}, \qquad \text{poiesis} = \text{threshold region between named states}.
\end{equation}
rose は named thing である。bud も named thing である。だが、その between には innumerable な phases がある。それらの phases は、名前をつけるのが hard だという merely その reason によって unreal なのではない。contrary に、Vuong は、wonder はそこで live すると言う。threshold に関する Heidegger の language が help するのは、ordinary naming が skip しようとする question を ask するからである。exactly where において rose は rose になるのか。

これは merely textbook 的な poiesis account に対する crucial な improvement である。Vuong は Aristotle を historical decoration として invoke しているのではない。mimetic naming が give out し、living process が begin する exact place を mark するために Aristotle を use しているのである。

\subsection{Question \& Answer}

subject は in itself cliché なのか、それとも私たちがそれを estrange できないときに only cliché になるのか。

この lecture において、subject は guilty ではない。treatment が guilty なのである。subject が dead になるのは、stale な recognition を通してそれに approach するときだけである。したがって practical rule は次のようになる。
\begin{enumerate}
\item familiar object を ban しない。
\item その object に対する first available treatment を refuse する。
\item それを reposition し、re-time し、あるいは rename する。
\item recognition がふたたび seeing へ open するまで、それと stay する。
\end{enumerate}
estrangement は world の rejection ではない。habit から world を recover することなのである。

\section{syntax、haunting、そして reading の pattern logic}

ひとたび perception が reopened されると、lecture は new puzzle へ shift する。もはや how we see だけを ask していては enough ではない。その seeing がどうやって reader の中で survive するのかも ask しなければならない。Vuong は、その problem を bluntly に state する。彼は、reader を hook する方法より、reader と stay する方法に interested なのだ、と。host は、この point を readerly terms へ translate し続ける。そして result は lecture の clearest formulations の一つになる。

Vuong の ratio は rough だが useful である。
\begin{equation}
80\% \text{ looking and thinking}, \qquad 20\% \text{ syntax}.
\end{equation}
彼は immediately に、listener が likely に抱く misunderstanding を correct する。final twenty percent は ``everything'' なのだ、と。なぜなら syntax は downloading mechanism だからである。したがって lecture の working model は
\begin{equation}
\text{perception} + \text{syntactic shaping} \Longrightarrow \text{readerly stain / haunting}.
\end{equation}
となる。

Siken の example は、これが practice においてどう work するかを show する。stars は ``little boats rowed out too far'' になる。tenor は十分 familiar である。
\begin{align}
\text{tenor} &= \text{stars},\\
\text{vehicle} &= \text{boats rowed out too far}.
\end{align}
だが correspondence が、whole symbolic regime を change する。ordinarily grand で mythic である stars が、intimate で lonely で late で vulnerable になるのだ。この line は、merely fresh な comparison を offer しているのではない。それは night sky の emotional scale 全体を change しているのである。

Ben Lerner は、それと同じ demand の harsher な pedagogical version を supply する。student の line は ``decent'' だ。だが Google は、hundreds of thousands がすでにそれを書いていることを reveal する。この blow は memorable である。lecture の opening criterion を sharpen するからだ。line が perception を置き去りにしたままなら、competence では enough ではないのである。

ここから lecture は naturally に haunting へ turn する。Browning の Meeting at Night は、それが paraphrase に reduce できないがゆえに、Vuong の中に stay している。Good Will Hunting の park-bench scene も同じである。love や experience に関する flat な sentence で concept は deliver できたかもしれない。だが、その scene は、それが vulnerable で image-rich な syntax へ pass したときに powerful になる。この stretch において、host の role は crucial である。何が happened したのかを、彼は plainer language で restate し続ける。ten seconds earlier には、それは recognition だった。いまやそれは seeing の one way になっているのだ。

argument はそれから、sentence から art in general へ broaden する。
\begin{equation}
\text{pattern} \Longrightarrow \text{either satisfy expectation or deny expectation}.
\end{equation}
sentence は linear technology であるがゆえに、time において unfold する。music もそうだ。film editing も。DJ set も。skate video も。ひとたび sequence が exist すれば、expectation が exist する。そしてひとたび expectation が exist すれば、maker は二つの fundamental resources を gain する。
\begin{equation}
\text{satisfy/deny expectation} \Longrightarrow \text{delight} + \text{surprise} + \text{estrangement}.
\end{equation}

\paragraph{Pattern Rule.}
\begin{enumerate}
\item linear art は、time において unfold することによって expectation を create する。
\item maker はその expectation を satisfy することも、hold off することもできる。
\item delay は、withheld な completion を audience が sense するため、attention を intensify する。
\item release、あるいは release の refusal が、delight と surprise を create する。
\item その point で、私たちはもはや merely content を following しているのではない。pattern 自体の intelligence を feeling しているのである。
\end{enumerate}

だからこそ Vuong は、literature、street-ball deception、skate parts、DJ timing のあいだをこれほど easily に move できるのだ。analogy は ornamental ではない。structural なのである。

\subsection{Question \& Answer}

sentence は、merely hook するのではなく、どうすれば reader と stay できるのか。

hook は eye を moment のあいだ capture する。haunting はその moment のあとにも remain する。Vuong の account において、sentence が stay するのは、その syntactic shape が perception の pressure に adequate であるときである。line は、merely concept を communicate するのでは enough ではない。rhythm、delay、pressure をもって reader の inside に enter しなければならない。だからこそ abstraction は、body を与えられたときに only memorable になるのだ。

\section{language の laboratories としての poetry と nature writing}

readerly after-effect から、lecture は again broaden する。こんどは artistic formation へ向かう。question は、where can a sentence still take risks? である。Vuong の answer は、poetry と nature writing は、most contemporary prose が partially surrender してしまった freedoms を preserve している、ということだ。

first claim は direct である。
\begin{equation}
\text{assignment removed} \Longrightarrow \text{language becomes available for experiment}.
\end{equation}
second claim は、それに follow する。
\begin{equation}
\text{poetry},\ \text{nature writing} \Longrightarrow \text{highest local freedom for the sentence}.
\end{equation}
poetry が obvious な laboratory であるのは、one need not maintain plot、explanatory clarity、character-management を every instant において maintain しなくてよいからだ。nature writing が parallel な laboratory になるのは、more subtle な reason による。pure mimesis がそこで collapse するからである。meadow はすでに見ている。mud がどんなものかも already know している。ならば、available な sight を merely duplicate する sentence をなぜ read すべきなのか。

J.A. Baker の passage が、その question に demonstration によって答える。Baker は plain な weather と marsh description から begin し、それから proliferating な mud-syntax に enter する。そこでは material がもはや merely named されない。mud は pressure、texture、danger、condition、psychic field になる。by the end には、Baker の interiority が object の中へ leached into しているのである。
\begin{equation}
\text{mud description} + \text{released interiority} \Longrightarrow \text{mimesis broken}.
\end{equation}
これが exactly、notes が preserve すべき movement である。Baker は world を abandon しているのではない。subjective force によって、world が page の上で how it can look を change させているのである。

host はここで another essential term を contribute する。fun である。この word が matter するのは、そうでなければ lecture が all theory and denunciation になってしまうからだ。many blues を含む crayon box の delight が、artistic perception の model になる。shades、textures、uses を multiply するとき、私たちは merely more finely classify しているのではない。wonder を reopen しているのである。だからこそ lecture は、strain なく Baker の mud から childlike attention へ move できるのだ。

Shklovsky-Tolstoy material は、この point で slightly different register において return する。Tolstoy は ``birch'' とは言わない。いわば、white luminous trunk をもつ large curly-headed tree と言う。point は descriptive excess ではない。child を again believe させることなのだ。definition は、object を too early に close するとき、imagination の enemy になりうる。

だからこそ lecture の late stretch において Eduardo C. Corral の moss simile は so important なのである。example の just before の transcript は unstable だが、technical claim は secure である。comparison は image-to-image ではなく、behavior-to-behavior において work する。
\begin{align}
\text{applause} \not\sim \text{moss by image}, \qquad \text{applause} \sim \text{moss by behavior},\\
\text{behavioral correspondence} \Longrightarrow \text{apparent acceleration of moss growth}.
\end{align}

\paragraph{Worked Example: Moss and Applause.}
\begin{enumerate}
\item moss が statically に何に resemble するかを ask するのではない。
\item moss がどう behave するのか、あるいはその behavior をどう newly perceive できるのかを ask する。
\item ``applause'' は quickened され、nebulous で、collective な motion を moss に lend する。
\item transferred な behavior が、moss を our eyes の前で move するように seem させる。
\item したがってこの simile は、merely decorate するのではなく、target の temporal experience を change するのである。
\end{enumerate}

Corral が short book に nine years を spent したという anecdote は、Babel とともに始まった loop を close するがゆえに matter する。strong metaphor とは instant cleverness ではない。prolonged looking の result なのである。

\subsection{Question \& Answer}

なぜ poetry と nature writing は、most contemporary prose より better に sentence の freedom を preserve するのか。

両方の forms が、immediate な assignment の pressure を loosen するからである。poetry は language itself に focus できる。nature writing は、visible object が already known されているため、mere report としては survive できない。したがって両者は、sentence を poiesis、estrangement、play へ force する。laboratory とは simply、world に more than one use を discover することが writer に permitted されている space のことなのである。

\section{publishing time、reader time、そして daringness の ethics}

late lecture はいま、local craft から larger systems へ shift する。house style、pedagogical cynicism、Hollywood conservatism、``reader in the Midwest''、そして Lotman の two-time model が、one question の parts になっていく。なぜ system は originality を principle の上では celebrate しながら、practice においてはそれを sameness へ discipline するのか。

Vuong はすでに、sentence が how standardized されるかを示している。ここで彼は、institutions がその standardization を actively preserve する仕方を示す。The New Yorker の example が matter するのは、それが simple denunciation ではないからである。house style は clarity を teach する。同時に、それは brand-recognition を produce する。writer は there に publish し、その piece が still idea においては自分のものだが、cadence においては no longer wholly theirs であることに気づくかもしれない。point は、house style が evil だということではない。clarity、efficacy、brand-expectation は never neutral だ、ということなのである。

``reader in the Midwest'' に関する anecdote は、その cynicism を sharpen する。editor は、そのような reader が baroque prose を handle できるかと ask する。Vuong がその question に hear するのは democratic care ではなく condescension である。ordinary reader も nervous system、memory、reading life を持っている、と彼は insist する。この anecdote が Uri Lotman への transition を prepare する。
\begin{equation}
\text{reading} = (\text{synchronic line},\text{diachronic line}).
\end{equation}
synchronic reading は season、publicity cycle、market year、institutional moment の inside で occur する。diachronic reading は、time を through した longer accumulation of reading に対して occur する。

これによって lecture の strongest institutional contrast が produce される。
\begin{equation}
\text{system judges synchronically}, \qquad \text{reader judges diachronically}.
\end{equation}
publishing は seasons の中で work する。review culture も seasons の中で work する。critics もしばしばそうである。readers はそうではない。reader は、today の prize-listed novel が arrive する long before から、Shakespeare、Melville、Baldwin、Annie Dillard、Chaucer を read してきたかもしれない。したがって synchronic success のために standardized された book は、diachronic hands に fall したとき、very different な test に meet することになる。

host の Rotten Tomatoes analogy は、ここで perfect である。Lotman の theory を familiar な cultural split へ translate してくれるからだ。critics は more institutional present の中に embedded されている。audiences は less regulated な accumulation of comparison を伴って arrive する。lecture の point は、audiences は pure で critics は corrupt だということではない。two are operating on different clocks だということなのだ。

Vuong はそれから、publication-time から language-time へ return する。definition は constrict することもあれば、expand することもある。good dictionaries や etymologies は、words を close するのではなく reopen する。passion の example は decisive である。modern usage が enthusiasm または intensity として hear するものは、older history として suffering を carry している。その history を recover した瞬間、word は altered される。これが Wittgenstein への bridge である。
\begin{equation}
\text{meaning of a word} = \text{its use}.
\end{equation}
use が definition を change する。dictionary は私たちに catch up する。この lesson が students にとって especially important なのは、official rule への intimidation が standardization へ戻る another path だからである。

ここから lecture は、again widen する。lexicon から culture へ。
\begin{equation}
\text{innovation on the margins} \Longrightarrow \text{capture by the center} \Longrightarrow \text{homogenized product}.
\end{equation}
margin は mobile である。center は capture する。Lotman の concentric model of culture は、innovation が why so often recognized center の outside で begin し、そのあと flattened commodity として return してくるのかを explain してくれる。同じことが words、styles、scenes、forms にも hold する。

host はそのあと、most practical late question of all を ask する。writer は actual に何をすべきなのか。Vuong の answer は another technical rule ではない。ethic である。daringness と disobedience である。daringness とは、risk を wager し、failure の risk を take する willingness である。disobedience とは、line-keeping が praise を bring する merely その reason で line に step back することを refuse することである。skateboarding analogy が matter するのは、daringness に felt form を与えるからだ。one throws oneself toward a trick without certainty of landing. failure は accidental byproduct ではない。practice itself の part なのである。

lecture は最後に、craft を beyond し、philosophical risk へ push する。language は Vuong の life を made してきた。それは livelihood、access、enlargement of perception を彼に与えてきた。だが same language は tyranny によって capture されうる。newspapers や radio は、authoritarian control の first targets の among である。したがって closing paradox は intact のまま remain しなければならない。
\begin{align}
\text{language} &\Longrightarrow \text{power} + \text{livelihood} + \text{enlargement of perception},\\
\text{language} &\Longrightarrow \text{capture} + \text{tyranny} + \text{homogenization}.
\end{align}
Thomas Thistlewood の library、Nazi officers の cultivated brutality、その他の historical examples が、lecture が sentimental になるのを prevent する。同時に、Harriet Beecher Stowe は ledger の other side に remain する。literature は sometimes history を alter する。Vuong は、この両方の truths を at once に insist するのである。

\subsection{Question \& Answer}

なぜ culture は originality を ask しながら conformity を reward するのか。そして writer はそれについて何を supposed to do すべきなのか。

institutions が recognizable product を need するからである。standardization は、edit しやすく、brand しやすく、circulate しやすく、praise しやすい。writer の answer は、この lecture においては purity ではなく practice である。
\begin{enumerate}
\item surprise の original matrix を protect すること。
\item seasonally だけではなく diachronically に read すること。
\item official intimidation より living use を trust すること。
\item daringness と disobedience を together に cultivate すること。
\item language が powerful であることを accept しつつ、それが innocent だと pretend しないこと。
\end{enumerate}
lecture は、この path が rewarded されるという guarantee を offer しない。これより anything less では sentence は live するには too timid になってしまう、という claim を offer するだけである。

\section{まとめ}

lecture は steady widening によって unfold する。metaphor についての question は workshop についての question になる。workshop についての question は sentence の history についての question になる。その history は estrangement、threshold、syntax、pattern、poetry、nature writing、publication-time、reader-time、そして finally language それ自体の moral ambiguity へ open する。each step において argument は again local example へ contract する。Babel の sunset、botanist の method、displaced された rose、Siken の stars、Baker の mud、Corral の moss、reader in the Midwest である。

end に remain するのは、stylistic maximalism の doctrine ではない。more exact な standard of attention である。私たちは longer に look し、mere recognizability を distrust し、sentence を consciousness の carrier として hear し、formal daring は artistic であると同時に ethical な wager でもあることを remember するよう ask されている。lecture は、それが end すべき場所で end する。resolution においてではなく、language が can do することと、it cannot promise することとについての heightened seriousness において。